\chapter*{Abstract}
\addcontentsline{toc}{chapter}{Abstract}

Digital Mohalla is a neighbourhood resource-sharing platform for lending, borrowing, and donation. Its complex computing challenge is not merely providing item listings: it is deciding whether a physical handover has sufficient trustworthy evidence when interaction history is sparse, indirect, and potentially subject to manipulation, while strictly preserving location privacy and escalating uncertain cases to human moderation.

This report documents the FYP-1 mid-assessment progress and implementation status frozen at 13 October 2026. Prior to implementation, a complete design baseline (Part A) covering use cases, sequence diagrams, ERD, state machines, and ADRs was integrated. In implementation (Part B), Slice A provides an event-sourced, request-to-return borrowing lifecycle: users authenticate via argon2-hashed credentials and JWT tokens, borrowing requests are validated against temporal constraints, stored locality eligibility is verified, and transactions proceed through accept, decline, handover, and return confirmations backed by an append-only, tamper-evident PostgreSQL hash-chain event log. 

Slice B completes the core trust and risk assurance engine: outcome extraction from historical transactions, recency-weighted Beta estimation with 90\% credible intervals, contextual similarity weighting, trust-edge projection, bounded $k$-hop graph search with exponential decay and noisy-OR aggregation, context-aware decision policies (Rules R1--R9), structured explanations, real FastAPI integration, synthetic scenario suites, and a three-layer mobile/web lender evidence screen. Concurrency controls prevent overlapping accepted bookings using database-level PostgreSQL exclusion constraints. Slice C remains in progress: collusive manipulation detection and dispute/moderation workflows remain to be completed. Multi-tiered test suites (unit, property-based Hypothesis, integration, and scenario tests) validate reproducibility.
